% =====================================================================
\section{Comparison with existing quantum codes}
\label{sec:cmp}
% =====================================================================

\subsection{The terms of the comparison}

The comparison is with the tables of \cite{GrasslCodetables}, retrieved
2026-11-04 (page last changed 10.06.2024). Two limitations must be stated
before any number is quoted, and we state them where they bite.

\begin{enumerate}[label=\textnormal{(\arabic*)},leftmargin=*]
\item \textbf{The tables cover stabilizer codes only.} There is no
  entanglement-assisted table: the corresponding endpoint returns HTTP 404.
  Every entanglement-assisted code of this paper is therefore
  \emph{unbenchmarked}. This is a substantial blind spot, since the
  entanglement-assisted route is where the largest distances of this paper
  occur.
\item \textbf{Only $q=4$ and $q=9$ are tabulated.} For $q=8$ the site returns
  ``currently not available for q=8''. Hence the $(8,\cdot,\cdot)$ codes of this
  paper are unbenchmarked as well. (A third point of convention: the site prints
  the \emph{prime} $p$ in its code labels, so that $q=4$ appears as
  $[[N,K]]_2$ and $q=9$ as $[[N,K]]_3$, while the construction block says
  ``over GF($p^2$)''. An earlier phase of this project misread this and
  concluded, incorrectly, that $q=9$ was not covered; we record that as an
  erratum in the supplementary material and the corrected comparison is the one
  below.)
\end{enumerate}

\subsection{The comparison table}

\begin{table}[t]
\centering
\caption{Every stabilizer code ($c=0$) produced by the framework, compared
with \cite{GrasslCodetables}. ``Best known'' is the lower .. upper bound
pair for $[[N,K]]_q$.}
\label{tab:comparison}
\begin{tabular}{@{}lllc@{}}
\toprule
code & family $(q,t,n)$ & route & best known $[[N,K]]_q$ / gap\\
\midrule
$[[16,12,2]]_4$ & $(4,1,4)$ & R2-SO & $2..2$ --- meets\\
$[[16,12,2]]_4$ & $(4,1,4)$ & R2-DC & $2..2$ --- meets\\
$[[20,16,2]]_4$ & $(4,1,5)$ & R2-SO & $2..2$ --- meets\\
$[[20,16,2]]_4$ & $(4,1,5)$ & R2-DC & $2..2$ --- meets\\
$[[24,20,2]]_4$ & $(4,1,6)$ & R2-DC & $2..2$ --- meets\\[2pt]
$[[20,12,2]]_4$ & $(4,1,5)$ & R2-SO & $3..3$ --- 1 below\\
$[[20,12,2]]_4$ & $(4,1,5)$ & R2-DC & $3..3$ --- 1 below\\
$[[24,16,2]]_4$ & $(4,1,6)$ & R2-DC & $3..3$ --- 1 below\\[2pt]
$[[16,4,3]]_4$  & $(4,1,4)$ & R2-SO & $5..5$ --- 2 below\\
$[[16,4,3]]_4$  & $(4,1,4)$ & R2-DC & $5..5$ --- 2 below\\
$[[20,4,4]]_4$  & $(4,1,5)$ & R2-SO & $6..6$ --- 2 below\\
$[[20,4,4]]_4$  & $(4,1,5)$ & R2-DC & $6..6$ --- 2 below\\[2pt]
$[[16,8,3]]_9$  & $(9,1,4)$ & R2-DC & $4..4$ --- 1 below\\
$[[24,8,4]]_9$  & $(9,1,6)$ & R2-SO & $6..8$ --- 2 below (best known 6)\\
$[[24,12,3]]_9$ & $(9,1,6)$ & R2-DC & $5..6$ --- 2 below (best known 5)\\
\bottomrule
\end{tabular}
\end{table}

\subsection{Reading of the comparison}

Every stabilizer code constructed here is at best \emph{equal} to the best known
distance, and then only with $d=2$, where equality is uninformative. Every code
with $d\ge3$ is one or two below the best known value. \textbf{No record
parameter is claimed}, and none is implied: the codes are presented for their
origin, not for their position in the tables.

That said, the comparison does say something. The framework produces, at
$N=24$ over $\F_9$, a stabilizer code of distance $4$ where the best known is
$6$ (and where the upper bound is $8$): the construction is within a factor of
$1.5$ of the best known code at that length and dimension, and it does so from a
chain of entirely algebraic operations over a non-commutative ring. More
importantly, the \emph{entanglement-assisted} codes --- where this framework's
distances are largest ($d=10$ at $N=20$) --- cannot be compared at all, and the
absence of an entanglement-assisted table is a limitation of the literature
rather than of the construction.

\subsection{Novelty}

We do not claim that matrix rings, Galois hulls, Gray maps, quasi-cyclic images,
stabilizer codes or entanglement-assisted codes are new; each is attributed in
Section~\ref{sec:lit}. The novelty claim is the following, specific and proved.

\begin{enumerate}[label=\textnormal{(\arabic*)},leftmargin=*]
\item \textbf{The duality.} Over the non-commutative ring $R=M_2(A_t)$, the
  Galois twist must be $\theta_s=\sigma_s\circ(\cdot)^{\mathsf T}$, an
  \emph{anti}-automorphism. Over a commutative ring this degenerates to
  $\sigma_s$ and there is nothing to prove; here it is what makes the dual of a
  \emph{right} module a right module with the correct size
  (Theorems~\ref{thm:duality} and~\ref{lem:failure}).
\item \textbf{Hulls without principality, and the refutation.} The entire hull
  theory is developed for arbitrary right ideals, and the natural conjecture
  that Galois hulls are principal is \textbf{false}
  (Theorem~\ref{thm:nonprincipal}), with the phenomenon confined by the
  structural condition $t\ge2$ and $p\mid n$. This has no commutative analogue.
\item \textbf{The exact induced form.} The coefficient expansion transports the
  Galois pairing to $u^{[s]}(I_{4n}\otimes H_t)v^{\mathsf T}$
  (Theorem~\ref{thm:gram}) --- a $\widehat\sigma_s$-Hermitian form with a
  Hankel Gram matrix, not the Euclidean or plain Hermitian form unless $t=1$.
  Duality and hulls are preserved including on the non-principal hulls
  (Theorems~\ref{thm:grayduality} and~\ref{thm:grayhull}).
\item \textbf{Self-orthogonality without reflexivity.}
  Theorem~\ref{thm:hullso} is proved by a one-sided argument, so the whole
  construction is available in the general Galois regime.
\item \textbf{Two new bounds.} The coset-distance characterisation
  (Theorem~\ref{thm:cosetdistance}) and the coset Singleton bound
  (Theorem~\ref{thm:cosetsingleton}), the latter being neither the
  entanglement-assisted Singleton bound nor implied by it
  (Proposition~\ref{prop:newsbound}).
\item \textbf{Quantum codes from non-principal hulls.} The codes
  $[[24,6,2;2]]_9$ and $[[32,8,2;4]]_4$ of Table~\ref{tab:headline} come from
  hulls with $\MU_{\Bq}=3$. Their parameters are weak; their existence is the
  point, since it shows that no principality hypothesis is hidden anywhere in
  the chain.
\end{enumerate}

Relative to the closest prior work \cite{ZhangKongZheng2026QuantumGaloisHulls}
--- Galois hulls, a Gray map, and quantum codes over a finite non-chain ring ---
the differences are: their ring is \emph{commutative}, so no anti-automorphism
arises and no principality question can be asked; their Gray image has length
$tn$, whereas ours has length $4tn$ (Theorem~\ref{thm:nphys}), because
$\dim_{\F_q}M_2(A_t)=4t$; they determine hull \emph{dimensions}, whereas we
construct codes from hulls that are not principal and classify when
$\MU_{\Bq}\ge3$ can occur; and we prove the two bounds of item (5).

\subsection{On the earlier tuple $[[24,4,5;12]]_4$}

An earlier phase of this project reported $[[24,4,5;12]]_4$, obtained at
$(q,t,n)=(4,2,3)$ from a $300$-code sample. Exhaustive enumeration of that
family ($35\,937$ codes) shows that its best entanglement-assisted distance is
$6$, attained by $[[24,4,6;20]]_4$ ($k=20$, $h_E=0$, $c=20$, $K=4$): the
earlier value was a sampled lower bound, not an optimum. The tuple
$[[24,4,5;12]]_4$ is not reproduced by the exhaustive search and is dominated
(by $[[20,4,6;12]]_4$). We report this as a change of value obtained by a
stronger method, not as a retraction: nothing in the earlier derivation was
wrong, it was simply incomplete.
